% ============================================================
% MODELO DE RELATÓRIO TÉCNICO
% Uso: auditoria, prestação de contas, diagnósticos
% ============================================================
\documentclass[12pt,a4paper]{article}

\usepackage[utf8]{inputenc}
\usepackage[T1]{fontenc}
\usepackage{lmodern}
\usepackage[brazil]{babel}
\usepackage{geometry}
\usepackage{parskip}
\usepackage{setspace}
\usepackage{fancyhdr}
\usepackage{lastpage}
\usepackage{xcolor}
\usepackage{graphicx}
\usepackage{tabularx}
\usepackage{booktabs}
\usepackage{hyperref}
\usepackage{tocloft}
\usepackage{enumitem}
\usepackage{datetime2}

\geometry{top=3cm,bottom=2cm,left=3cm,right=2cm}
\setstretch{1.5}

\pagestyle{fancy}
\fancyhf{}
\fancyhead[L]{\small\textbf{Relatório Técnico}}
\fancyhead[R]{\small\today}
\fancyfoot[C]{\thepage/\pageref{LastPage}}
\renewcommand{\headrulewidth}{0.4pt}
\renewcommand{\footrulewidth}{0.4pt}

%% ── DADOS ──
\newcommand{\municipio}{Inajá}
\newcommand{\uf}{PR}
\newcommand{\tituloRelatorio}{Título do Relatório}
\newcommand{\responsavel}{Nome do Responsável}
\newcommand{\periodo}{Janeiro a Dezembro de 2026}

\begin{document}

%% ── CAPA ──
\begin{titlepage}
  \begin{center}
    \vspace*{3cm}
    {\large\textbf{PREFEITURA MUNICIPAL DE \MakeUppercase{\municipio} -- \uf}}\\[0.5cm]
    \rule{\textwidth}{0.4pt}\\[1cm]
    {\LARGE\textbf{\tituloRelatorio}}\\[1.5cm]
    {\large\textbf{Período:} \periodo}\\[0.5cm]
    {\large\textbf{Responsável:} \responsavel}\\[2cm]
    \vfill
    {\large \municipio/\uf, \today}
  \end{center}
\end{titlepage}

%% ── SUMÁRIO ──
\tableofcontents
\newpage

%% ── SEÇÕES ──
\section{Introdução}

Contexto e objetivo do relatório.

\section{Metodologia}

Procedimentos adotados para a análise.

\section{Constatacões}

Resultados e evidências encontradas.

\subsection{Subseção 1}

Detalhamento.

\subsection{Subseção 2}

Detalhamento.

\section{Conclusão e Recomendações}

Conclusão final e recomendações.

\section{Assinatura}

\vspace{1cm}
\rule{6cm}{0.4pt}\\
\textbf{\responsavel}

\end{document}
